\subsection{Quantum K-theory of a point}%
\label{sub:Quantum K-theory of a point}

One natural question is to study higher genus quantum K-theory. We will begin with the simplest target, which is a point. One motivation for studying higher-genus would be to try to prove integrality of the Klemm-Pandharipande definition of GV invariants for Calabi-Yau fourfolds or the Pandharipande-Zinger definition of GV invariants for Calabi-Yau fivefolds (in this situation, we only need to consider genus \(1\)). One reason we need to consider the quantum K-theory of a point is because in \(I\bar{M}_{1,1}(X, \beta)\), there are some components where an entire curve of genus \(1\) is contracted to a point.

\begin{quest}
    What is the K-theoretic Witten-Kontsevich theorem?\footnote{According to You-Cheng, the difficulty lies everywhere.}
\end{quest}

We will first recall the cohomological Witten's conjecture. We will begin with the operator
\[ L = \partial_x^2 + u(x) \]
and deform it to
\[ \tilde{L} = \partial_x^2 + u(x, T_3, T_5, \ldots), \]
where we should think of \(x = T_1\) and \(u(x) = u(x, 0, 0,\ldots)\). Formally, write
\[ L^{1/2} = \partial_x + \sum_{k > 0} a_{-k}(x) \partial_x^{-k} \]
using the property
\[ \partial^n \circ f = \sum_{m = 0}^{\infty} \binom{n}{m} f^{(m)} \partial^{n-m} \]
for all integers \(n \in \Z\). In particular, we construct the (2-)KdV hierarchy by the equation
\[ \pdv{\tilde{L}}{T_{2k+1}} = \ab[(\tilde{L}^{\frac{2k+1}{2}})_{\geq 0}, \tilde{L}] \]
for all \(k \geq 0\). We define a \textit{\(\tau\)-function} to be a function \(\tau\) such that 
\[ 2 \pdv[2]{}{x}\log \tau(x, T_3, T_5, \ldots) = u(x, T_3, T_5, \ldots). \]
Now perform the above process using \(u = 2x\). Define the function
\[ F(t_0, t_1, \ldots) \coloneqq \sum_{n \geq 0} \sum_{k_0, \ldots, k_n} \ab<\tau_0^{k_0} \cdots \tau_n^{k_n}> \frac{t_0^{k_0} \cdots t_n^{k_n}}{k_0! \cdots k_n!}. \]
Note that because the genus is determined entirely by the number of insertions and the degree of the insertions, the notation
\[ \ab<\tau_0^{k_0} \cdots \tau_n^{k_n}> \]
uniquely determines a descendant GW invariant of a point.

\begin{thm}[Kontsevich {\cite{Kontsevich1992}}]
    \(F\) is a \(\tau\)-function of the KdV hierarchy constructed starting with \(L = \partial_x^2 + 2x\) under the change of variables \(t_i = (2i+1)!! T_{2i+1}\).
\end{thm}

One application of this theorem is to consider virtual localization. For example, suppose we consider \(\bar{M}_{g,0}(X, d[C])\) where \(C\) is a unique rational curve in a Calabi-Yau threefold with normal bundle \(N = \mc{O}(-1)\oplus \mc{O}(-1)\). Then in this setting we can reduce to the local model and compute
\begin{align*}
    \int_{[\bar{M}_{g,0}(\P^1, d)]^{\vir}} c_{\mr{top}}(R^1 \pi_* N) = \sum_{m \vdash d} \frac{(-1)^{d-L(m)}}{\abs{\Aut(m)}} \int_{\bar{M}_{g, L(m)}} \frac{1 + \cdots + c_g(\mc{H})}{\prod_i (1-m_i \psi_i)}
\end{align*}
via virtual localization. If we would like a similar result in K-theory, then we will need to consider terms of the form
\[ \chi\ab(\bar{M}_{g, L(m)}/\Aut(m), \frac{\lambda_{-1}(\mc{H})}{1-m_i L_i}), \]
so we will need to consider the permutation-equivariant quantum K-theory of a point.

We now proceed with a literature review. First, Lee~\cite{Lee1997} computed the formula
\[ \ab<\frac{1}{1-q_1 L}, \ldots, \frac{1}{1-q_n L}>_{0,n}^{\mr{pt}, K} = \ab(\prod_{i=1}^n \frac{1}{1-q_i}) \ab(1 + \sum_{i=1}^n \frac{q_i}{1-q_i})^{n-3}. \]
Later, Lee--Qu~\cite{LeeQu2014} gave a way to compute
\[ \ab<\frac{1}{1-q_1 L}, \ldots, \frac{1}{1-q_n L}>_{1,n}^{\mr{pt}, K} \]
in principle. Givental~\cite{Givental2017K1} computed the formula
\[ \sum_{n \geq 2} \ab<\frac{1}{1-q L}, 1, \ldots, 1>_{0,1+n}^{S_n} = (1-q) \exp\ab(\sum_{k > 0} \frac{\Psi^k}{k} \ab(1 + \frac{q}{1-q})). \]
Finally, Milanov~\cite{Milanov2025}, in a preprint first posted in 2024, computed 
\begin{align*}
    \sum_{n \geq 0}& \ab<\frac{1}{1-q_1 L}, \ldots, \frac{1}{1-q_r L}, 1, \ldots, 1>_{0,r+n}^{S_n} \\
    &= \ab(\prod_{i=1}^{r} \frac{1}{1-q_i})\ab(1 + \sum_{i=1}^r \frac{q_i}{1-q_i}) \exp\ab(\sum_{k > 0} \frac{\Psi^k}{k} \ab(1 + \sum_{i=1}^r \frac{q_i}{1-q_i})). 
\end{align*}

The following pluri-Hodge recursions and the new genus-one calculations based on them are due to You-Cheng Chou and were presented in this lecture.

Let \(s \in \Z\) and define the pluri-Hodge sheaf \(\mc{H}_s \coloneqq R^0 \pi_* (\omega_{\pi}^s)\). Considering the map \(\bar{M}_{g,n+1} \xrightarrow{\pi} \bar{M}_{g,n}\), we will denote the upstairs cotangent lines by \(L_i\) and the downstairs ones by \(\ell_i\).

\begin{lem}
    For all integers \(s\) and \(d_i\), we have
    \[ \pi_*\ab(\omega_{\pi}^s\ab(\sum_{i=1}^n d_i D_i)) = \mc{H}_s-\mc{H}_{1-s}^* + \sum_{i=1}^n \sum_{k=1}^{d_i} \ell_i^{-k+s}. \]
\end{lem}

\begin{rmk}
    Note that if \(f(q) \in \Z[q]\), then \(\exp \sum_k \frac{\Psi^k}{k} (f(q)) \in \Z \ps{q}\). The reason is that this operation sends \(f = \sum_{i=1}^{n} a_i q^i\) to \(\prod_{i=1}^n \frac{1}{(1-q^i)^{a_i}}\).
\end{rmk}

\begin{proof}
    When all \(d_i\) are zero, we see that
    \[ R^0 \pi_* \omega_{\pi}^s - R^1 \pi_* \omega_{\pi}^s = \mc{H}_s - \mc{H}_{1-s}^* \]
    by Serre duality. The next step is to show that
    \[ \pi_* \omega_{\pi}^s ((r+1)D_i) = \pi_* \omega_{\pi}^s(r D_i) + \ell_i^{s-r-1}. \]
    Here, we first consider the exact sequence
    \[ 0 \to \mc{O} (-D_i) \to \mc{O} \to \mc{O}_{D_i} \to 0 \]
    and twist by \(\omega_{\pi}^s\) to obtain
    \[ 0 \to \omega_{\pi}^s(r D_i) \to \omega_{\pi}^s ((r+1)D_i) \to \omega_{\pi}^s ((r+1)D_i) \otimes \mc{O}_{D_i} \to 0 \]
    and then compute the pushforward.
\end{proof}

\begin{lem}
    Let \(f_i\) be rational functions, for example \(\frac{1}{1-q L}\). Then
    \begin{align*}
        \ab<f_1(L), \ldots, f_n(L), L_{n+1}^r>_{g, n+1} ={}& \sum_{i=1}^n \ab<f_1(L), \ldots, D(f_i f_{n+1}), \ldots, f_n(L)>_{g, n} \\
        &+ \ab<f_1(L), \ldots, f_n(L) \mid \mc{H}_r - \mc{H}_{1-r}^*>_{g,n},
    \end{align*}
    where \(f_{n+1} = L_{n+1}^r\) and \((Df)(x) \coloneqq \frac{f(x)-f(1)}{x-1}\).
\end{lem}

To prove this, we need to consider the pushforward
\begin{align*}
    \pi_* (L_1^{d_1} \cdots L_{n+1}^{d_{n+1}}) &= \pi_* \ab(\pi^*(\ell_1^{d_1} \otimes \cdots \otimes \ell_n^{d_n}) \otimes \omega_{\pi}^{d_{n+1}}\ab(\sum_i (d_i+d_{n+1})D_i)).
\end{align*}
Here, we use the fact that \(L_i = (\pi^* \ell_i)(D_i)\) and \(L_{n+1} = \omega_{\pi}(D_1 + \cdots + D_n)\). Applying the previous lemma, we see that the pushforward becomes
\[ (\ell_1^{d_1} \cdots \ell_n^{d_n})\ab(\mc{H}_{d_{n+1}} - \mc{H}_{1-d_{n+1}}^* + \sum_{i=1}^n \sum_{k=1}^{d_i + d_{n+1}} \ell_i^{-k+d_{n+1}}). \]
Now, Lee's formula computes the correlators
\[ \ab<(L-1)^{k_1}, \ldots, (L-1)^{k_n}>_{0,n} = \frac{(n-3)!}{\prod_{i=1}^n k_i! \ab(n-3-\sum_i k_i)!} \]
via the expansion
\[ \frac{1}{1-qL} = \frac{1}{1-q} \sum_{k \geq 0} \frac{q^k(L-1)^k}{(1-q)^k}. \]
Note here that there is no Hodge bundle in genus \(0\), so in fact the lemma can be used to obtain Lee's formula.

\begin{lem}
    We have
    \[ \mc{H}_r - \mc{H}_{1-r}^* = \pi^* (\mc{H}_r - \mc{H}_{1-r}^*) - \sum_{i=1}^n (D^2 (\pi^* \ell_i)^r) \mc{O}_{D_i}. \]
\end{lem}

The lemmas collectively imply that we can reduce to the case of a single marking at the cost of making the insertions very complicated and introducing pluri-Hodge terms everywhere. In genus \(1\), life becomes easier because on \(\bar{M}_{1,1}\), we have 
\[ \mc{H}_r = \mc{O}(r) = L_1^r. \]
In this case, recall the classical formula
\[ \chi\ab(\bar{M}_{1,1}, \frac{1}{1-qL}) = \frac{1}{(1-q^4)(1-q^6)}. \]
Unfortunately, computing
\[ \chi\ab(\bar{M}_{1,n}; \frac{1}{1-q_1 L} \cdots \frac{1}{1-q_n L}) \]
is too hard, so we can try to compute
\[ \chi\ab(\bar{M}_{1,1}; \frac{1}{1-q_1 L_1} \cdots \frac{1}{1-q_n L_1})  \]
which is known, but it is unclear how to write a good formula for it. The reason we would consider this is the operation
\[ \Z\ps{q} \xrightarrow{[]_r} \Z\ps{q_1, \ldots, q_r} \qquad q^k \to h_k(q_1, \ldots, q_r) \coloneqq \sum_{\ell_1 + \cdots + \ell_r = k} q_1^{\ell_1} \cdots q_r^{\ell_r}. \]
Unfortunately, You-Cheng does not know how to write
\[ \chi\ab(\bar{M}_{1,1}; \frac{1}{1-q_1 L_1} \cdots \frac{1}{1-q_n L_1}) = \ab[\frac{1}{(1-q^4)(1-q^6)}]_n \]
in a nice form. In principle we can write formulae like
\[ [f(q)]_2 \coloneqq \frac{q_1 f(q_1) - q_2 f(q_2)}{q_1 - q_2}, \]
but this does not seem to help much.

\subsection{Permutation-equivariant theory of a point}%
\label{sub:Permutation-equivariant theory of a point}

We can compute \(\chi\ab(\bar{M}_{1,1}, \frac{1}{1-q L})\) in another way. Consider the involution on an elliptic curve, which gives a map
\[ \bar{M}_{1,1} \to \bar{M}_{0,1+3}/S_3. \]
Then, Givental~\cite{Givental2017K1} computes
\[ \chi\ab([\bar{M}_{0,1+n-1}/S_{n-1}], \frac{1}{1-qL}) = \prod_{i=2}^{n-1} \frac{1}{1-q^i}. \]
In particular, we obtain
\[ \frac{1}{(1-q^2)(1-q^3)}, \]
and then because the map is a \(\mu_2\)-gerbe (or a \(2:1\)-cover on objects), we obtain the desired formula.

Consider \(\mc{H}_{g,1} \subset \bar{M}_{g,1}\) to be the hyperelliptic locus. Note that the marking lies on a Weierstrass point. Note that the fixed points of the hyperelliptic involution go to \(2g+1\) points on \(\P^1\), and therefore one would want to use the calculation
\[ \chi\ab(\bar{M}_{0,1+2g+1}/S_{2g+1}, \frac{1}{1-q L}) = \prod_{i=2}^{2g+1} \frac{1}{1-q^i} \]
to compute the quantity
\(\chi\ab(\mc{H}_{g,1}, \frac{1}{1-qL})\). However, because of the way that the curves upstairs degenerate, we will need to consider the Hassett moduli space~\cite{Hassett2003} where the last \(2g+1\) markings have weight \(\frac{1}{2}\).
 






